\section{Conclusion}

Finding quantum-vulnerable cryptography is the first step every migration
plan requires, and it cannot be done from source code alone. \tool{} scans
five layers and reconciles them. It finds \data{rqone-recall-ratio} as many
labeled assets as the best open-source alternative, and it shows a federal web
where \data{rqtwo-ep-share} of TLS endpoints already negotiate hybrid key
exchange, largely because of where they are hosted, while every leaf
certificate remains classical. \data{rqthree-any-share} of the most critical open-source
code bases we scanned carry quantum-vulnerable cryptography, most often in configuration and key
material. Hybrid key exchange costs a kilobyte-scale ClientHello and, at the
median, at most \data{rqfour-netem-extra-max} per handshake in our
measurements, so its cost is unlikely to explain where it is missing.

\paragraph{Availability.} \tool{}, the corpus, the harness, the measurement
scripts, the raw data and this paper's generator are open source at
\url{https://github.com/emmanueladutwum123/pqscope}.
